% Living draft of the research paper. Move to the official ICML template once the 2027 call is posted.
% Numbers in tables come from the run logs: python agent/results_table.py (see paper/README.md).
\documentclass[11pt]{article}
\usepackage[margin=1in]{geometry}
\usepackage{amsmath}
\usepackage{booktabs}
\usepackage{xcolor}
\usepackage{hyperref}
\newcommand{\todo}[1]{\textcolor{red}{[TODO: #1]}}
\newcommand{\status}[1]{\textcolor{blue}{[Status: #1]}}

\title{Agentic In-Process Redesign of 3D-Printed Lattices\\ \large Living draft, not for submission}
\author{\todo{authors and order, to agree with Yichen. ICML review is anonymous.}}
\date{Last results: experiment E6, 2026-11-05}

\begin{document}
\maketitle

\begin{abstract}
\status{placeholder until the main result exists.}
When a 3D print develops a defect partway through, the usual choices are to continue and accept a weaker part, or to restart. We study a third option: pause the print, and have an AI agent redesign the part that has not been printed yet, so that the finished part still meets its original mechanical requirement. So far, we have built the test cases, an automatic checker, and comparison methods, and run a first AI agent that can only change wall thickness. In that restricted setting, a calculator tool greatly improved the agent, but it did not beat a simple hand-written rule, and it was about level with a non-AI search using the same number of calculator uses. All mechanical results so far are formula estimates, not physical tests.
\end{abstract}

\section{Introduction}
\status{outline.}
A printed part may already contain a measurable defect while material remains to be printed. Earlier layers cannot be changed, but later layers can. The research question is whether a constrained change to the remaining geometry can recover the part's original compression response, and whether an AI agent adds anything over simpler ways of choosing that change.

Three claims must be established separately:
\begin{enumerate}
  \item \textbf{Recovery is physically possible:} a valid change improves the response compared with continuing the original design.
  \item \textbf{Manufacturing information matters:} designs that use the measured defect beat generic ones.
  \item \textbf{The AI adds useful capability:} with the same tools, data and options, the agent beats a non-AI method.
\end{enumerate}
The current experiments address only part of claim 3, and only on formula estimates.

\todo{Explain why this belongs at ICML: what is the machine-learning contribution? Agree with Yichen.}

\section{Related work}
\status{to write from the proposal's literature review.}
\todo{In-process monitoring and correction in additive manufacturing; AI models for printer control and CAD editing; re-entrant lattice mechanics.}
Our agent loop follows the action and observation pattern of ReAct \cite{yao2023react}. Stiffness and strength scaling follow \cite{gibson1997cellular}. A planned experiment tests writing the prompt as code \cite{puerto2024code}.

\section{Problem setup}
\status{describes the current implementation.}
\paragraph{Specimen.} A re-entrant honeycomb, 50 $\times$ 50 $\times$ 10 mm, printed along Z in 0.2 mm layers with 1.2 mm walls and loaded along Y. Hexagon and triangle designs exist but have not been used in agent experiments.

\paragraph{Cases.} There are 12 cases: four defect types (under-extrusion, over-extrusion, layer shift, warpage), at medium severity, starting at 10\%, 50\% or 80\% of the height. Each case assumes the defect is noticed 1 mm later and the print pauses there. Everything below the pause height is fixed.

\paragraph{A fix.} One wall thickness $t_i$ for each remaining 1 mm band, applied to all walls in that band.

\paragraph{Rules.} A fix is allowed if all of these hold: $0.8 \le t_i \le 2.0$ mm; $t_i - t_{i-1} \le 0.4$ mm, where $t_0$ is the last printed wall thickness; and plastic used is at most the plastic the original design would use above the pause height, plus 1\%.

\section{Scoring}
\status{a stand-in; expected to change.}
Let $r_k$ be the wall area of layer $k$ divided by the wall area of a healthy layer, excluding the end plates. Printed layers are measured by slicing the 3D file, and proposed layers are computed from the design. Treating the $N = 50$ layers as springs side by side, and using scaling for bending-dominated cells \cite{gibson1997cellular}:
\begin{equation}
  \text{Stiffness} = \frac{1}{N}\sum_{k=1}^{N} r_k^{3}, \qquad
  \text{Peak force} = \frac{1}{N}\sum_{k=1}^{N} r_k^{2}.
\end{equation}
Both are relative to a healthy part, so 100\% is the target, and error is the distance from 100\% in either direction. For the triangle design, which is stretch-dominated, both exponents are 1.

\paragraph{Known limits.} These formulas ignore which walls changed. They cannot see layer shift, because a shifted layer keeps its area, and they barely register warpage. They ignore bonding between layers, and they reward alternating thick and thin bands more than a physical test may. The proposal's real target is the measured force--displacement curve.

\section{Methods compared}
\begin{description}
  \item[Do nothing.] Keep the original design.
  \item[Simple rule.] A fixed recipe. It measures how much stiffness the printed part has lost or gained, then thins all remaining bands if too stiff, or moves plastic into alternating thin and thick bands if too weak.
  \item[Search.] COBYLA, a standard non-AI optimizer, starting from the original design. It learns about a design only through the calculator, with a limit of 10 uses for a fair comparison. A run allowed 200 uses serves as a reference for what this program can reach.
  \item[Agent, no calculator.] An AI model receives the facts of the case and proposes a fix. A rejected fix returns with reasons, up to 3 attempts, then falls back to the original design.
  \item[Agent, calculator.] The same, plus up to 10 calculator uses before submitting. The calculator returns plastic used, the plastic limit, rule violations, and estimated stiffness. It does not return peak force.
\end{description}

\section{Results}
\status{preliminary.}
Experiment E6 ran both agents on all 12 cases, 3 times each, on Nemotron 3 Ultra with hidden reasoning turned off.

\begin{table}[h]
\centering
\begin{tabular}{lrrrr}
\toprule
& \multicolumn{2}{c}{All 12 cases} & \multicolumn{2}{c}{6 extrusion cases} \\
Method & Stiffness & Peak force & Stiffness & Peak force \\
\midrule
Do nothing & 3.05\% & 2.00\% & 5.86\% & 3.83\% \\
Simple rule & 0.79\% & 0.90\% & 1.34\% & 1.64\% \\
Search, 10 uses & 1.01\% & 0.82\% & 1.91\% & 1.56\% \\
Search, 200 uses & 0.41\% & 0.60\% & 0.79\% & 1.19\% \\
Agent, no calculator & 2.76\% & 1.84\% & 5.27\% & 3.52\% \\
Agent, calculator & 1.17\% & 0.96\% & 2.15\% & 1.80\% \\
\bottomrule
\end{tabular}

\caption{Average distance from a healthy part, lower is better. Only the 6 extrusion cases separate the methods, because the score barely registers layer shift and warpage. Search and rule results are deterministic; agent values average 3 runs per case.}
\end{table}

\begin{table}[h]
\centering
\begin{tabular}{lrrrrr}
\toprule
Case & Do nothing & Simple rule & Search (10) & Agent, no calc. & Agent, calc. \\
\midrule
under extrusion 10\% & 95.5\% & 100.3\% & 97.9\% & 95.1\% & 97.5\% \\
under extrusion 50\% & 95.5\% & 97.9\% & 97.5\% & 95.7\% & 95.8\% \\
under extrusion 80\% & 95.5\% & 95.5\% & 95.9\% & 95.5\% & 95.6\% \\
over extrusion 10\% & 107.2\% & 100.2\% & 101.1\% & 107.2\% & 100.1\% \\
over extrusion 50\% & 107.2\% & 99.6\% & 99.0\% & 107.2\% & 100.2\% \\
over extrusion 80\% & 107.2\% & 100.5\% & 100.5\% & 103.5\% & 101.5\% \\
layer shift 10\% & 100.0\% & 100.0\% & 100.0\% & 100.0\% & 100.0\% \\
layer shift 50\% & 100.0\% & 100.0\% & 100.0\% & 100.0\% & 100.0\% \\
layer shift 80\% & 100.0\% & 100.0\% & 100.0\% & 100.0\% & 100.0\% \\
warpage 10\% & 99.5\% & 99.5\% & 99.5\% & 99.5\% & 99.8\% \\
warpage 50\% & 99.5\% & 99.5\% & 99.9\% & 99.5\% & 100.0\% \\
warpage 80\% & 99.5\% & 99.5\% & 99.9\% & 99.5\% & 99.5\% \\
\bottomrule
\end{tabular}

\caption{Stiffness as a percentage of a healthy part, target 100\%. The number in each case name is the height where the defect starts.}
\end{table}

\paragraph{Observations.}
\begin{itemize}
  \item The calculator improved the agent, with no case-level losses against the agent without one.
  \item On the 6 extrusion cases, the agent with a calculator had 0 wins, 3 ties and 3 losses against the simple rule, and 2 wins, 2 ties and 2 losses against the search at 10 uses.
  \item On under-extrusion, the agent mostly tried thickening every band, hit the plastic limit, and returned to the original design in 5 of 9 runs. It rarely found the idea of moving plastic between bands.
  \item This outcome is consistent with the proposal's expectation that an ordinary optimizer may be the best method when only a few thickness settings can change.
\end{itemize}

\section{Limitations}
\begin{itemize}
  \item All mechanical values are formula estimates, with no physical tests or simulation.
  \item The calculator's stiffness estimate uses the same formula as the score, so the agent can aim directly at the main measure. Peak force is the only measure it cannot see.
  \item Defects are modeled in the design files, and are not produced by a printer.
  \item One lattice shape, one severity, and three runs per case.
  \item Hidden reasoning was off, which may understate the model. On Gemini 3.6 Flash with reasoning on, an earlier run of the agent without a calculator averaged a lower error than it did here.
\end{itemize}

\section{What the paper still needs}
\begin{enumerate}
  \item \todo{A scoring method independent of the agent's calculator, that can see layer shift and warpage.}
  \item \todo{The full redesign setting, where the agent changes more than thickness, with the search given the same options.}
  \item \todo{More shapes, severities and heights, with enough runs to report uncertainty.}
  \item \todo{Some physical validation, or a clear statement of why there is none. Decide with Yichen.}
  \item \todo{Move to the official ICML template and check the page limit when the 2027 call is posted.}
\end{enumerate}

\bibliographystyle{plain}
\bibliography{references}
\end{document}
